\section{Конечная разность — дискретная производная}

Метод конечных разностей — дискретный аналог дифференциального и интегрального исчисления, где
производную заменяет конечная разность, а интеграл — суммирование.

\begin{definition}
Для функции дискретного аргумента $f(x)$ (или последовательности $a_k=f(k)$, $x=1,2,\dots$)
\emph{(правой) конечной разностью} называют
\[
  \Delta f(x)=f(x+1)-f(x).
\]
Это аналог производной $f'(x)=\lim_{h\to0}\frac{f(x+h)-f(x)}{h}$ при шаге $h=1$.
\end{definition}

\begin{example}
$\Delta 1=0$; $\ \Delta x=(x+1)-x=1$; $\ \Delta x^2=(x+1)^2-x^2=2x+1$.
\end{example}

\begin{property}[свойства оператора $\Delta$]\leavevmode
\begin{enumerate}[label=\arabic*),leftmargin=*]
  \item $\Delta(c\,f)=c\,\Delta f$, $\ c\in\R$;
  \item $\Delta(f+g)=\Delta f+\Delta g$;
  \item $\Delta(f\cdot g)=\Delta f(x)\cdot g(x)+\Delta g(x)\cdot f(x+1)$ (дискретное правило Лейбница).
\end{enumerate}
\end{property}

\section{Факториальные степени}

Степенная функция $x^m$ неудобна тем, что $\Delta x^m\ne mx^{m-1}$. Её роль в дискретном анализе
играют факториальные степени.

\begin{definition}
Для $m\in\N$ \emph{убывающая факториальная степень}
\[
  x^{\underline{m}}=\underbrace{x(x-1)(x-2)\cdots(x-m+1)}_{m\text{ множителей}},\qquad
  x^{\underline 0}=1 .
\]
Для $m\in\N$ \emph{отрицательная} факториальная степень
$\displaystyle x^{\underline{-m}}=\frac{1}{(x+1)(x+2)\cdots(x+m)}$.
\end{definition}

\begin{theorem}
Для любого целого $m$ справедливо $\Delta\bigl(x^{\underline m}\bigr)=m\,x^{\underline{m-1}}$ —
точный аналог $\bigl(x^m\bigr)'=mx^{m-1}$.
\end{theorem}

\begin{proof}
Для $m\in\N$:
\[
  \Delta x^{\underline m}=(x+1)x(x-1)\cdots(x-m+2)-x(x-1)\cdots(x-m+1)
   =x^{\underline{m-1}}\bigl[(x+1)-(x-m+1)\bigr]=m\,x^{\underline{m-1}}.
\]
Случай $m\le0$ проверяется аналогично прямой подстановкой.
\end{proof}

\section{Неопределённая сумма и дискретная формула Ньютона–Лейбница}

\begin{definition}
\emph{Неопределённой суммой} (дискретной первообразной) называют оператор $\Sigma$, обратный к
$\Delta$:
\[
  \sum f(x)\,\delta x=g(x)+C\quad\iff\quad \Delta g(x)=f(x)
\]
(обозначение Д.~Кнута). Аналогично $\displaystyle\sum x^{\underline m}\,\delta x
=\frac{x^{\underline{m+1}}}{m+1}+C$ для $m\ne-1$.
\end{definition}

\begin{theorem}[дискретная формула Ньютона–Лейбница]
Для целых $a\le b$
\[
  \sum_{x=a}^{b-1}\Delta f(x)=f(x)\Big|_a^b=f(b)-f(a).
\]
Эквивалентно (определённая сумма Кнута):
$\displaystyle\sum_{x=a}^{b-1}f(x)=\sum f(x)\,\delta x\,\Big|_a^b$.
\end{theorem}

\begin{proof}
Сумма телескопическая:
\[
  \sum_{x=a}^{b-1}\bigl(f(x+1)-f(x)\bigr)=\bigl(f(a+1)-f(a)\bigr)+\dots+\bigl(f(b)-f(b-1)\bigr)
   =f(b)-f(a).
\]
\end{proof}

\subsection{Таблица дискретных первообразных}

\begin{center}
\renewcommand{\arraystretch}{1.5}
\begin{tabular}{ll}
$\Delta(x^{\underline{m}})=m\,x^{\underline{m-1}}$ & $\displaystyle\sum x^{\underline m}\delta x
  =\frac{x^{\underline{m+1}}}{m+1}+C\ (m\ne-1)$\\
$\Delta(2^{x})=2^{x}$ & $\displaystyle\sum 2^{x}\delta x=2^{x}+C$\\
$\Delta(c^{x})=c^{x}(c-1)$ & $\displaystyle\sum c^{x}\delta x=\frac{c^{x}}{c-1}+C\ (c\ne1)$\\
$\Delta(H_n)=\dfrac{1}{n+1}$ & $\displaystyle\sum \frac{1}{n+1}\,\delta n=H_n+C$\\
$\Delta(F_n)=F_{n-1}$ & $\displaystyle\sum F_n\,\delta n=F_{n+1}+C$\\
\end{tabular}
\end{center}
(здесь $H_n$ — гармонические числа, $F_n$ — числа Фибоначчи; роль $\sum c^x$ — точный аналог
$\int a^x\,dx=\tfrac{a^x}{\ln a}$).

\begin{example}
$\displaystyle\sum_{x=1}^{10}\frac{1}{x(x+1)(x+2)}=\sum_{1}^{11}x^{\underline{-3}}\,\delta x
=\frac{x^{\underline{-2}}}{-2}\Big|_1^{11}
=\frac12\Bigl(\frac{1}{1\cdot2}-\frac{1}{11\cdot12}\Bigr)=\frac{65}{264}.$
\end{example}

\section{Суммирование многочленов}

Многочлен раскладывают по факториальным степеням и суммируют почленно. Коэффициенты разложения
дают \emph{дискретная формула Тейлора} и \emph{таблица конечных разностей}.

\begin{theorem}[дискретная формула Тейлора]
Для многочлена $P_m$
\[
  P_m(n)=\sum_{k=0}^{m}\frac{\Delta^k P_m(0)}{k!}\,n^{\underline k},
\]
где $\Delta^0 f=f$, а $\Delta^k$ — $k$-кратная разность (аналог формулы Маклорена с заменой
$f^{(k)}\to\Delta^k$, $n^k\to n^{\underline k}$).
\end{theorem}

\begin{example}[сумма квадратов]
Для $f(n)=n^2$ таблица разностей в нуле даёт $\Delta^0 f(0)=0$, $\Delta f(0)=1$,
$\Delta^2 f(0)=2$, поэтому $n^2=n^{\underline1}+n^{\underline2}$. Тогда
\[
  \sum_{n=a}^{b-1}n^2=\sum_a^b\bigl(n^{\underline1}+n^{\underline2}\bigr)
   =\Bigl(\frac{n^{\underline2}}{2}+\frac{n^{\underline3}}{3}\Bigr)\Big|_a^b .
\]
\end{example}

\begin{example}[сумма кубов]
$n^3=n^{\underline1}+3n^{\underline2}+n^{\underline3}$ (из таблицы разностей: $\Delta^k n^3(0)$
равны $0,1,6,6$). Поэтому
\[
  \sum_{n=1}^{N}n^3=\Bigl(\frac{n^{\underline2}}{2}+n^{\underline3}+\frac{n^{\underline4}}{4}\Bigr)
   \Big|_1^{N+1}=\frac{(N+1)^2 N^2}{4}=\Bigl(\frac{N(N+1)}{2}\Bigr)^2 .
\]
\end{example}

\section{Суммирование по частям (формула Абеля)}

\begin{theorem}[суммирование по частям]
\[
  \sum u(x)\,\Delta v(x)=u(x)v(x)-\sum \Delta u(x)\cdot v(x+1),
\]
или в определённой форме
$\displaystyle\sum_{x=a}^{b-1}u(x)\,\Delta v(x)=u(x)v(x)\Big|_a^b-\sum_{x=a}^{b-1}\Delta u(x)\,
v(x+1)$.
\end{theorem}

\begin{example}
$\displaystyle\sum_{x=2}^{N}2^{x}\,x(x-1)$. Берём $u=x^{\underline2}$, $\Delta v=2^x$ (т.е.
$v=2^x$):
\[
  \sum_{2}^{N+1}x^{\underline2}\,\Delta(2^x)=x^{\underline2}2^{x}\Big|_2^{N+1}
   -\sum_{2}^{N+1}2x^{\underline1}\,2^{x+1},
\]
и повторное применение формулы доводит сумму до замкнутого вида — дискретный аналог
$\int x^2 e^x\,dx$ двукратным интегрированием по частям.
\end{example}
